% Manuscript skeleton designed for an eight page limit.
% Replace every \fillin{...}, \blankcell, and \blankwide after analysis.
% Detailed protocol and secondary analyses belong in sections/appendix.tex.

\section{Introduction}

Large language models (LLMs) can suggest codes, organize codebooks, and apply qualitative categories across large collections of text \citep{dai2023llm,wang2026centaurta,matveyenko2026muse}. Researchers often review these outputs, edit suggestions, explain disagreements, and guide later analysis. Most evaluations still treat their judgments as one general source of human feedback, even when the task calls for different kinds of knowledge.

A qualitative methods expert may recognize that a topic label has been presented as an interpretive theme. A domain expert may notice that an analysis has lost a cultural distinction or omitted a clinically important exception. A researcher without either specialization may still identify an unsupported quotation or an overly broad claim. A single average score conceals these differences and provides little guidance about how to use limited expert time.

Evaluations of LLM supported clinical thematic analysis remain fragmented across expert review and automatic similarity measures \citep{yi2025position}. The usual measures do not resolve this problem. More than one code space may be defensible \citep{chen2026opencodes}, agreement between humans and LLMs may not predict blind expert preference \citep{liu2026agreement}, and model confidence can be poorly calibrated \citep{li2026confidence}. Audits have also found missing participant evidence, summaries presented as quotations, and failures that depend on local context \citep{parkington2025human,bang2026hybrid,subbiah2026story}. Provenance identifies the source of a passage. It does not show that the passage supports the interpretation or represents the range of views in the corpus.

The planned experiment treats review as an \emph{allocation problem}: who should examine a qualitative interpretation and its supporting evidence? \textbf{WarrantRoute} will collect judgments from researchers without specialist expertise, qualitative methods experts, and domain experts. The study will include natural model outputs and controlled variants with known defects. A simple policy will send each item to no specialist, a qualitative methods expert, a domain expert, or both. The policy will use information available before specialist review, and all comparisons will receive the same amount of expert time. We will freeze the policy before testing it on a new domain and model family. Figure~\ref{fig:warrantroute_framework} summarizes the framework.

\begin{figure*}[t]
    \centering
    \includegraphics[width=\textwidth]{figures/warrantroute_framework.pdf}
    \caption{Overview of the planned WarrantRoute evaluation. An interpretable policy will route source-grounded theme and evidence packets to researchers alone, a qualitative methods expert, a domain expert, or both within a fixed time budget. The routing comparison uses CaCHe and Dreaddit; KODIS, CaSiNo, and AMI remain supplementary, unassigned corpus lanes.}
    \label{fig:warrantroute_framework}
\end{figure*}

The planned study is designed to contribute:

\begin{enumerate}
    \item a source-grounded benchmark that isolates five validity problems without assuming a single correct thematic interpretation;
    \item a human-evaluation design that adapts evidential credibility and plural-reference coverage from prior work and adds study-specific, source-grounded operationalizations of voice-and-boundary preservation and scope calibration; and
    \item a planned held-out domain test that compares several allocation rules at equal cost and measures whether selected feedback repairs a defect without creating a new one.
\end{enumerate}
The present manuscript reports the protocol and a synthetic measurement-pipeline qualification, not a real-corpus human-evaluation result.

\section{Related Work}

\paragraph{Human feedback in qualitative analysis.}
Human-in-the-loop systems let analysts revise codes and rationales \citep{dai2023llm}, pair-code and discuss disagreements \citep{gao2024collabcoder}, convert repeated feedback into persistent instructions \citep{wang2026centaurta}, and retain positionality and revision histories \citep{ye2026reflexis}. Multi-agent systems extend this workflow: Auto-TA automates source-to-theme generation \citep{yi2025autota}, SFT-TA uses supervised fine-tuned agents \citep{yi2025sftta}, and TAMA incorporates cardiac expert guidance \citep{xu2026tama}. These studies establish scalable collaboration, but their evaluations emphasize alignment, usability, control, or aggregate output quality rather than reviewer specialization.

\paragraph{Review allocation and evaluator dependence.}
Selective-review methods route items using confidence and ensemble diversity \citep{zhao2025confidence} or rationale disagreement \citep{tajik2026disagreement}, although confidence can be miscalibrated or anchor-sensitive \citep{li2026confidence}. Evaluator identity also changes conclusions: blind expert verification can diverge from agreement metrics \citep{liu2026agreement}, and specialists who read source material can differ from nonspecialists who see isolated themes \citep{moskalewicz2026phenomenology}. Adjacent biomedical work evaluates generated protocol pseudocode against a generated reference \citep{yi2024protomed}, but not the source warrant of qualitative interpretations. Prior work therefore studies allocation signals and evaluator perspective separately, not which expert role a source-linked packet requires under matched expert time.

\paragraph{Source grounding and artifact quality.}
Provenance systems preserve links among quotations, codes, themes, and revisions \citep{ye2026reflexis,yi2026provenance}, yet models can omit human-identified excerpts, substitute summaries for quotations, or produce paraphrased, truncated, and hallucinated evidence \citep{parkington2025human,bang2026hybrid}. Computational thematic-analysis evaluations have operationalized credibility through human acceptance and semantic alignment and, later, binary LLM judgments \citep{qiao2025thematiclm,wang2026centaurta}; OpenCodes instead measures coverage relative to a pooled code space \citep{chen2026opencodes}. These constructs are complementary: cited evidence may support a claim that omits dissent, while broad coverage does not ensure appropriate scope. WarrantRoute therefore adapts evidential credibility and plural-reference coverage, adds packet-level operationalizations of voice-and-boundary preservation and scope calibration, and reports all four separately. The open comparison is whether methods and domain expertise detect complementary source-grounded failures, and whether expertise-aware routing improves detection over expertise-agnostic allocation at the same cost.

\section{Task and Study Design}

\subsection{Evidence packets and controlled variants}

The planned unit of analysis is a \textbf{theme and evidence packet}. Each future real-data packet will contain a qualitative claim, a short explanation, and exact excerpts only from an access-controlled, pseudonymized analysis copy after the required two-person privacy and permissions review, with project-pseudonymous speaker or document identifiers, local context, and a summary of source coverage. Raters will judge whether the evidence supports the claim and preserves important differences, counterevidence, and source voice. The task is a bounded source-warrant and representation assessment. It will not assess a complete reflexive thematic analysis or assume that one thematic account is universally correct \citep{braun2021can,braun2021quality}.

The study will sample natural LLM outputs and construct variants vetted by experts. Each variant will change one property frozen in the preregistration:

\begin{enumerate}
    \item \textbf{unsupported evidence}: a topically related excerpt does not warrant the claim;
    \item \textbf{source concentration}: wording implies breadth while evidence comes disproportionately from one participant or document;
    \item \textbf{counterevidence loss}: a rare, negative, or boundary case is removed;
    \item \textbf{contextual flattening}: a consequential cultural or domain specific qualification is erased; or
    \item \textbf{unsupported abstraction}: a local observation is generalized beyond its evidence.
\end{enumerate}

An independent construction panel will check the original packet, the intended defect, and the absence of a second manipulation. Pilot items will not enter the main analysis, and each rater will see at most one version of a base packet.

\subsection{Artifact-quality measures}

We will preregister four artifact-quality measures, separately from the routing and repair outcomes below. Two adapt measures used in prior computational thematic analysis. \textbf{Evidential credibility} asks whether the core interpretive meaning is directly supported by the cited evidence in local context. Unlike prior binary LLM judgments \citep{qiao2025thematiclm,wang2026centaurta}, the planned primary measure is a blinded human rating. \textbf{Plural-reference coverage} is the percentage of concepts in a frozen reference space pooled from multiple independent analysts that are recovered by at least one generated code or theme \citep{chen2026opencodes}. The pool will preserve alternative concepts rather than designate one analyst's themes as truth, and blinded experts will adjudicate many-to-many conceptual matches.

We propose two study-specific, packet-level operationalizations. \textbf{Voice-and-boundary preservation} assesses whether consequential differences, dissent, counterevidence, less-frequent or boundary cases, and contextual qualifications survive synthesis. \textbf{Scope calibration} assesses whether claim breadth, strength, polarity, and causal language are proportionate to the supplied evidence. Raters will score these two measures and evidential credibility from 1 to 5. The analysis will retain the ordinal scores and report ratings of 4 or 5 as descriptively adequate; ``cannot judge'' will be reported separately. Plural-reference coverage will be calculated only on complete natural-output inventories, separately for codes and themes, and will not be treated as a 1--5 rating. We will interpret the four measures only after deterministic quotation, attribution, assignment, and provenance checks. Because coverage is recall-like, we will report generated-concept volume and human-audited unsupported or redundant concepts so that verbosity cannot masquerade as quality. We will not combine the four measures into a composite. Full anchors and matching rules appear in Appendix~\ref{sec:artifact_metrics}.

\subsection{Datasets and leakage controls}

The project maintains five candidate corpus lanes that differ in domain,
interaction structure, and provenance. Dreaddit's official training split is
provisionally assigned to development, its official test split to a secondary
in-domain audit, and the 11 CaCHe focus groups to the
planned held-out cross-domain confirmatory evaluation. KODIS is a received,
aggregate-audited private exploratory dialogue lane. CaSiNo and the AMI
scenario-meeting subset remain source-inventory candidates. KODIS, CaSiNo, and
AMI have no frozen experimental role, approved model or rater pathway, or
sampling eligibility. Here, ``held out'' describes the planned procedure, not
model-training provenance.

Table~\ref{tab:dataset_statistics} reports each lane at its current preparation
stage. Dreaddit's candidate pool applies duplicate and exact-date cluster
quarantines; Dreaddit and CaCHe remain pending manual privacy review and project
authorization. KODIS counts exclude 775 human--AI dyads and nine human--human
rows without transcripts, and its exploratory parser retains repeated-role
turns. CaSiNo counts exclude 2,378 structured negotiation events, while AMI
counts transcribed speech segments and words. All lanes are
English-language or translated into English. The heterogeneous units are not
common experimental denominators, and unassigned rows are not sampling
eligible. Definitions, provenance, access conditions, and calculations appear
in Appendix~\ref{sec:corpora_controls}.

\begin{table*}[t]
\centering
\scriptsize
\setlength{\tabcolsep}{2pt}
\renewcommand{\arraystretch}{0.92}
\begin{tabularx}{\textwidth}{@{}l>{\raggedright\arraybackslash}Xl>{\raggedright\arraybackslash}Xl@{}}
\toprule
Dataset & Context & Language & Inventory & Stage \\
\midrule
Dreaddit \citep{turcan2019dreaddit} & Social-media stress posts & English & \shortstack[l]{2,929 posts; 3,553 segments;\\3,516 candidate; 301,279 words} & \shortstack[l]{Processed;\\review pending} \\
CaCHe \citep{mehta2024cache2} & Sexual and reproductive health focus groups & Translated English & \shortstack[l]{11 groups; 3,944 turns;\\3,076 eligible; 98,735 words} & \shortstack[l]{Processed;\\review pending} \\
KODIS \citep{hale2025kodis} & Customer-service dispute dialogues & English & \shortstack[l]{2,085 human dyads; 2,076 dialogues;\\27,121 messages; 568,333 words} & \shortstack[l]{Exploratory;\\review pending} \\
CaSiNo \citep{chawla2021casino} & Campsite negotiation dialogues & English & \shortstack[l]{1,030 dialogues; 11,919 messages;\\228,675 words} & Source only \\
AMI scenario \citep{carletta2006ami} & Product-design meetings & Spoken English & \shortstack[l]{138 meetings; 69,258 segments;\\793,764 words} & Source only \\
\bottomrule
\end{tabularx}
\caption{Dataset statistics.}
\label{tab:dataset_statistics}
\end{table*}

Under the planned protocol, only Dreaddit training data will inform prompts,
packet rules, controlled variants, routing features, and thresholds. We will
freeze these choices before test generation, preserve complete source groups,
quarantine exact training duplicates that overlap Dreaddit test text, and
cluster CaCHe analyses by focus group and transcript-local speaker. KODIS,
CaSiNo, and AMI will not inform these choices while their experimental roles remain
unassigned. Packets will use the same context limits across adopted evaluation
datasets. We will retain prompts, decoding settings, run identifiers, source
spans, and revisions. Protected text will be processed only under approved
terms. The complete leakage and generation protocol appears in
Appendix~\ref{sec:corpora_controls}.

\subsection{Raters and annotation}

The planned study will recruit three groups. \textbf{Qualitative methods experts} will have demonstrated experience conducting qualitative analysis. \textbf{Domain experts} will have research, professional, or relevant lived experience in the source domain. \textbf{Researchers without specialist expertise} will be able to apply the study rubric but will not meet either expert definition for the assigned item. Dual expertise will be recorded as a separate status. Item construction, initial rating, policy development, and final adjudication will be handled by different people wherever staffing permits.

All raters will complete the same tutorial and review balanced subsets in a randomized incomplete-block design. For each packet they will rate evidential credibility, voice-and-boundary preservation, scope calibration, and confidence on scales from 1 to 5. The common quality anchors will be 1 (not adequate), 2 (major problems), 3 (partially adequate; material revision needed), 4 (mostly adequate; minor issues only), and 5 (fully adequate); confidence will use separate uncertainty anchors. They will also choose accept, revise, reject, or escalate; indicate whether they cannot judge the item; provide a short rationale and, if needed, request a particular kind of expertise. The system will record review time. ``Cannot judge'' will remain a separate response category, with no conversion to a midpoint or a quality failure. Qualitative methods experts and domain experts will submit independent ratings before adjudication. When their disagreement reflects defensible interpretations, the item will remain plural or ambiguous.

We will choose the final sample size after a separate pilot by simulating the crossed structure of raters and packets. The preregistration will state the number of unique raters, base packets, ratings per group, exclusions, and the smallest effect or noninferiority margin that the design can support.

\section{Routing by Expertise}

The planned routing method will use information available before specialist review: researchers' ratings of evidential credibility, voice-and-boundary preservation, and scope calibration, plus disagreement, confidence, abstention, the requested kind of expertise, response time, and nonsensitive provenance measures such as excerpt count, source concentration, and quotation exactness. Plural-reference coverage will be evaluation-only because it relies on the frozen analyst reference space. The router will not use manipulation labels, expert scores from the test set, protected participant attributes, or test outcomes.

We will fit an interpretable regularized model that chooses among no escalation, review by a qualitative methods expert, review by a domain expert, or review by both. We will select its thresholds on the development data to maximize recall of serious errors within a fixed budget measured in expert minutes. We will then compare five conditions:

\begin{enumerate}
    \item \textbf{researchers only}: no specialist review;
    \item \textbf{both expert groups}: a qualitative methods expert and a domain expert review every item;
    \item \textbf{random}: route the same budget at random;
    \item \textbf{uncertainty only}: route using confidence, disagreement, and abstention; and
    \item \textbf{WarrantRoute}: predict both whether escalation is needed and which role is relevant.
\end{enumerate}

Every condition will use the same observed expert time. The planned confirmatory outcomes on the CaCHe cross-domain evaluation are recall of serious errors, false acceptance, and expert minutes per correctly escalated error; the same outcomes will be reported on the Dreaddit in-domain audit as secondary results. These policy-level outcomes are distinct from the four artifact-quality measures. A secondary experiment will use one fixed revision model and compare no feedback, feedback from researchers, feedback selected by WarrantRoute, and feedback from both expert groups. A separate blinded panel will judge whether the revision fixes the target defect and whether it creates a new unsupported claim, omission, or distortion.

\section{Evaluation and Analysis}

The planned analysis addresses three questions. First, for which errors do qualitative methods experts or domain experts differ materially from researchers without specialist expertise? Second, can the initial ratings identify those cases within a fixed budget of expert time? Third, does the selected feedback repair defects in a new domain and model family?

For the first question, separate cumulative-link mixed-effects models will analyze evidential credibility, voice-and-boundary preservation, and scope calibration; logistic mixed-effects models will analyze decisions and defect detection. Models will include packet version, evaluator group, error family, domain, and their preregistered interactions, with random effects for rater and base packet. Plural-reference coverage will be estimated separately at the complete-run level, not from sampled packets or controlled twins. For the second question, we will compare policies on the cross-domain evaluation with cluster-bootstrap intervals and a preregistered noninferiority contrast against review by both expert groups. For the third question, a logistic mixed model will relate successful repair to feedback condition and error family. Dreaddit results will be secondary; the CaCHe analysis will become confirmatory only after its manifest and analysis plan are frozen and preregistered.

We will report effect sizes and uncertainty intervals. Ordinal Krippendorff's $\alpha$, Gwet's AC2, raw agreement, and confidence calibration will describe measurement behavior; they will not establish validity on their own \citep{mcdonald2019reliability}. We will freeze and preregister confirmatory contrasts, multiplicity, and the CaCHe packet manifest before generating model outputs for that corpus. Development and audit results, comparisons of natural and controlled items, and other robustness checks will appear in Appendix~\ref{sec:appendix_results}. Table~\ref{tab:artifact_reliability} is reserved for reliability, and Table~\ref{tab:transfer_ablation} for transfer and routing-feature ablations. Confidence and router calibration will be reported separately in the appendix.

\section{Results}

This section separates the completed synthetic measurement-pipeline
qualification from the planned real-corpus reporting shell. The
qualification used four fictional templates from an earlier corpus design,
including KODIS-shaped and CANDOR-shaped proxies. It tests whether the four
measures can be calculated reproducibly; it provides no evidence about
Dreaddit, CaCHe, KODIS, CaSiNo, AMI, or any other real corpus. Replace every later
blank from a single frozen analysis export and reconcile all denominators
before writing the final abstract. Report the CaCHe confirmatory results before
the Dreaddit audit and development performance.

\subsection{Synthetic measurement-pipeline qualification}

We applied the paper-aligned measures to 60 outputs generated from entirely fictional packets: five recent model variants, four corpus-shaped packets, and three independent runs per model and packet. Three blinded model judges independently scored evidential credibility, voice-and-boundary preservation, and scope calibration using the manuscript anchors. For this engineering pilot, the benchmark's frozen designed constructs instantiate the reference space; each judge matched every reference separately to generated codes and themes, and a two-of-three panel rule determined recovery. In Table~\ref{tab:synthetic_paper_metrics}, the ordinal columns report the percentage of 36 blinded ratings per model scored 4 or 5, and code and theme coverage report panel-majority recovery across 63 frozen reference opportunities per level and model. Brackets are 95\% packet-by-run cluster-bootstrap intervals, and integrity is deterministic. All judges and candidate models are from one provider, so these are measurement-pipeline results rather than real-corpus or human-evaluation results.

The pilot exposes a ceiling effect that would be hidden by reporting adequacy alone. GPT-5.6 Sol, GPT-5.5, and GPT-5.6 Terra received 100\% adequacy on all three ordinal measures, while their deterministic integrity pass rates were 10/12, 9/12, and 5/12. Code coverage ranged only from 93.7\% to 95.2\%, and theme coverage from 85.7\% to 93.7\%. We therefore retain the four constructs as the source-grounded artifact profile but do not use them alone for model selection: integrity failures are exclusionary, full ordinal distributions and controlled-defect sensitivity are reported, and coverage is accompanied by output-volume and unsupported-concept diagnostics. Under that gate-first rule, GPT-5.6 Sol remains the remediation candidate rather than a frozen paper model.

% Generated by aggregate_paper_metrics.py from validated blinded records.
\begin{table*}[t]
\centering
\scriptsize
\setlength{\tabcolsep}{2.3pt}
\begin{tabular}{lrrrrrrr}
\toprule
Model & \shortstack{Integrity\\pass} & \shortstack{Evidential\\adequacy} & \shortstack{Voice/boundary\\adequacy} & \shortstack{Scope\\adequacy} & \shortstack{Code\\coverage} & \shortstack{Theme\\coverage} & \shortstack{Serious\\error} \\
\midrule
GPT-5.6 Sol & 10/12 & 100.0 [100.0, 100.0] & 100.0 [100.0, 100.0] & 100.0 [100.0, 100.0] & 93.7 [89.1, 98.4] & 85.7 [78.5, 93.4] & 0.0 [0.0, 0.0] \\
GPT-5.4 & 9/12 & 100.0 [100.0, 100.0] & 97.2 [91.7, 100.0] & 100.0 [100.0, 100.0] & 93.7 [89.1, 98.4] & 93.7 [89.1, 98.4] & 5.6 [0.0, 13.9] \\
GPT-5.5 & 9/12 & 100.0 [100.0, 100.0] & 100.0 [100.0, 100.0] & 100.0 [100.0, 100.0] & 93.7 [89.1, 98.4] & 92.1 [87.3, 96.8] & 0.0 [0.0, 0.0] \\
GPT-5.6 Luna & 8/12 & 80.6 [69.4, 91.7] & 80.6 [69.4, 91.7] & 88.9 [80.6, 97.2] & 95.2 [90.9, 100.0] & 85.7 [79.7, 91.9] & 8.3 [0.0, 16.7] \\
GPT-5.6 Terra & 5/12 & 100.0 [100.0, 100.0] & 100.0 [100.0, 100.0] & 100.0 [100.0, 100.0] & 95.2 [90.9, 100.0] & 88.9 [84.1, 93.7] & 0.0 [0.0, 0.0] \\
\bottomrule
\end{tabular}
\caption{Synthetic measurement-pipeline qualification.}
\label{tab:synthetic_paper_metrics}
\end{table*}


\subsection{Pre-human corpus checkpoint}

All five real corpora are now present and deterministically inventoried.
Dreaddit and CaCHe have validated restricted derivatives; KODIS has a private
exploratory turn conversion; and CaSiNo and AMI have verified source
inventories but no adopted adapters. These preparation statistics populate
Table~\ref{tab:dataset_statistics}, not performance outcomes. The remaining
result cells require independent ratings, analyst reference construction and
concept matching, observed review time, or blinded repair judgments.

\subsection{Study accounting}

The final manuscript will report \fillin{N TOTAL RATERS} recruited raters and \fillin{N RATINGS} analyzed judgments over \fillin{N PACKETS} packets. \fillin{REPORT EXCLUSIONS, BALANCE, AND THE TEST SET}. No corpus-size statistic can fill these fields; Table~\ref{tab:study_summary} begins at the human-feedback boundary and is reserved for study counts by experimental role.

\begin{table}[t]
\centering
\scriptsize
\setlength{\tabcolsep}{3pt}
\begin{tabularx}{\columnwidth}{>{\raggedright\arraybackslash}Xccc}
\toprule
Quantity & Dev. & In domain & Cross domain \\
\midrule
Source participants/documents & \blankcell & \blankcell & \blankcell \\
Base packets / controlled twins & \blankcell & \blankcell & \blankcell \\
Researchers without specialist expertise & \blankcell & \blankcell & \blankcell \\
Qualitative methods experts & \blankcell & \blankcell & \blankcell \\
Domain experts & \blankcell & \blankcell & \blankcell \\
Independent ratings & \blankcell & \blankcell & \blankcell \\
Median review minutes & \blankcell & \blankcell & \blankcell \\
\bottomrule
\end{tabularx}
\caption{Planned reporting shell for study counts by experimental role. Raters and packets will be reported as unique counts.}
\label{tab:study_summary}
\end{table}

\subsection{Artifact-quality profile}

The final manuscript will report CaCHe credibility, voice-and-boundary, and scope-calibration adequacy rates, plus code and theme plural-reference coverage, only after the confirmatory analysis is frozen and run. It will report these measures with the integrity prerequisite, denominators, and uncertainty, and provide full ordinal distributions, cannot-judge rates, matching uncertainty, unmatched reference concepts, and generated-concept volume in Appendix~\ref{sec:artifact_metrics}.

\subsection{Where expertise changes judgments}

\fillin{SUMMARIZE THE INTERACTION BETWEEN EVALUATOR GROUP AND FAILURE FAMILY, INCLUDING UNCERTAINTY}. Table~\ref{tab:role_sensitivity} reports predicted probabilities of detecting a defect. \fillin{STATE WHETHER THE TWO EXPERT GROUPS PROVIDED COMPLEMENTARY, REDUNDANT, OR INCONCLUSIVE JUDGMENTS}.

\begin{table*}[t]
\centering
\small
\setlength{\tabcolsep}{4pt}
\begin{tabular}{lcccc}
\toprule
Failure family & Researchers & Methods experts & Domain experts & Largest contrast \\
\midrule
Unsupported evidence & \blankwide & \blankwide & \blankwide & \blankwide \\
Source concentration & \blankwide & \blankwide & \blankwide & \blankwide \\
Counterevidence loss & \blankwide & \blankwide & \blankwide & \blankwide \\
Contextual flattening & \blankwide & \blankwide & \blankwide & \blankwide \\
Unsupported abstraction & \blankwide & \blankwide & \blankwide & \blankwide \\
\bottomrule
\end{tabular}
\caption{Predicted probability of detecting each failure, with 95\% confidence intervals, by evaluator group.}
\label{tab:role_sensitivity}
\end{table*}

\subsection{Routing under matched expert time}

After the expert-minute budget is preregistered and the study is run, the final manuscript will report \fillin{PRIMARY TEST RESULT}, \fillin{EFFECT VERSUS UNCERTAINTY ROUTING WITH 95\% CONFIDENCE INTERVAL}, and \fillin{NONINFERIORITY RESULT VERSUS BOTH EXPERT GROUPS}. Table~\ref{tab:routing} is reserved for the five routing policies under the matched expert-time budget. Results must be reported by failure family if an average gain conceals misses on negative cases.

\begin{table*}[t]
\centering
\scriptsize
\setlength{\tabcolsep}{2.5pt}
\begin{tabular}{lccccc}
\toprule
Policy & Expert minutes & Serious error recall & False acceptance & Difference from both experts & Errors per expert hour \\
\midrule
Researchers only & 0 & \blankcell & \blankcell & \blankwide & Not applicable \\
Random & \blankcell & \blankcell & \blankcell & \blankwide & \blankcell \\
Uncertainty routing & \blankcell & \blankcell & \blankcell & \blankwide & \blankcell \\
WarrantRoute & \blankcell & \blankcell & \blankcell & \blankwide & \blankcell \\
Both expert groups & \blankcell & \blankcell & \blankcell & Reference & \blankcell \\
\bottomrule
\end{tabular}
\caption{Planned reporting shell for routing performance on the CaCHe confirmatory evaluation at matched budgets measured in expert minutes. Intervals will use a bootstrap clustered by base packet.}
\label{tab:routing}
\end{table*}

\subsection{Does routed feedback repair defects?}

The final manuscript will report \fillin{RATE AND CONFIDENCE INTERVAL} for packets repaired by WarrantRoute feedback without a new error, alongside \fillin{RESEARCHER FEEDBACK RATE} for researcher feedback and \fillin{BOTH EXPERT GROUPS RATE} for feedback from both expert groups. Table~\ref{tab:repair} is reserved for repair outcomes, new errors, credibility changes, and expert time. \fillin{REPORT PERFORMANCE RELATIVE TO EXPERT MINUTES AND THE MOST IMPORTANT FAILURE CASE}.

\begin{table}[t]
\centering
\scriptsize
\setlength{\tabcolsep}{2.8pt}
\begin{tabular}{lcccc}
\toprule
Feedback & Repair & New error & \shortstack{Credibility\\change} & \shortstack{Expert\\minutes} \\
\midrule
None & \blankwide & \blankwide & \blankwide & 0 \\
Researchers & \blankwide & \blankwide & \blankwide & 0 \\
WarrantRoute & \blankwide & \blankwide & \blankwide & \blankcell \\
Both expert groups & \blankwide & \blankwide & \blankwide & \blankcell \\
\bottomrule
\end{tabular}
\caption{Planned reporting shell for blinded revision. A repair will count as successful only if it fixes the targeted defect without introducing an unsupported claim.}
\label{tab:repair}
\end{table}

\section{Discussion}

\fillin{STATE WHICH EVALUATOR GROUP ADDED VALUE FOR WHICH FAILURE TYPES, WITH EFFECTS}. \fillin{STATE WHETHER THE TWO EXPERT GROUPS PROVIDED COMPLEMENTARY JUDGMENTS}. \fillin{REPORT THE PRACTICAL COST AND EXPERT-REVIEW COVERAGE CONSEQUENCE AT THE OBSERVED REVIEW TIMES}. Any resulting conclusions must be limited to the evidence-review task, failure families, and evaluator groups actually studied.

On the new domain and model family, \fillin{STATE THE TRANSFER RESULT AND ITS INTERPRETATION}. \fillin{STATE WHAT THE RESULT IMPLIES FOR ROUTING BEYOND GENERIC UNCERTAINTY}. A useful review queue should preserve the reason for escalation and allow a reviewer to request another form of expertise. Persistent disagreement remains visible for deliberation.

\section{Limitations}

Any resulting estimates will reflect the recruited reviewers, whose judgments may reproduce disciplinary, cultural, or institutional assumptions. Controlled variants will isolate individual failure mechanisms, but naturally occurring errors may be subtler or combine several defects. Repeated judgments from the same people may increase precision without increasing the diversity of expertise. Routing can also create false reassurance around items that are not escalated, so the design includes a random audit and failure-family-specific false-acceptance reporting. The task covers bounded evidence packets. A complete qualitative analysis also requires familiarization, memoing, reflexivity, and accountable interpretation. Generalization will be limited to English Reddit stress segments and English translations of focus groups from western Kenya, as well as \fillin{MODEL FAMILIES AND RECRUITMENT SAMPLE}. If the unassigned dialogue lanes are later adopted, KODIS and CaSiNo cannot support cross-dialogue participant counts or participant-disjoint splits from their public releases, while AMI scenario meetings are elicited role play and transcript-only processing omits prosody, timing, overlap, and multimodal evidence.

\section{Ethics Statement}

Real text model processing and rater exposure begin only after a documented
institutional determination and all source, model provider, rater access,
retention, and release gates are satisfied; source collection approvals do not
transfer. We treat all corpora as potentially identifying. Analysis copies are
minimized and restricted, protected text goes only to approved services, and
every excerpt shown to raters or released must pass a documented two-person
privacy and permissions review. Raters
participate voluntarily outside supervisory or evaluative relationships, with
identifiers stored separately. Public artifacts contain only aggregate,
synthetic, or cleared material. Routing supports research review, not clinical
or participant decisions. Details appear in Appendix~\ref{sec:detailed_ethics}.

\section{Conclusion}

WarrantRoute is designed to measure how reviewer expertise affects the detection and repair of unsupported qualitative interpretations under a fixed review budget. The current manuscript establishes the protocol and reports only the fictional-packet measurement qualification; \fillin{STATE THE PRIMARY REAL-CORPUS RESULT AND ITS PRACTICAL IMPLICATION ONLY AFTER THE FROZEN STUDY IS COMPLETE}.
